\chapter{Exception and Interrupt Handling}
\label{ch:exceptions}

\section{Overview}

An exception is an event that causes the program flow to be diverted to an address stored in a predefined vector.

Some are directly caused by an instruction, some are a side effect of an instruction, and some are asynchronously
caused by external hardware.

They are managed by the exception coprocessor, which takes control of the processor when an exception is pending.

\subsection{Instructions Used in This Chapter}

This chapter refers to several exception-coprocessor instructions before their full reference entries appear. Briefly:

\begin{description}
	\item[\mnemonic{EXCP}] Triggers a software exception (trap).
	\item[\mnemonic{RETE}] Returns from an exception handler, restoring the saved program counter, status register, and
	      exception level.
	\item[\mnemonic{WAIT}] Halts the CPU until an enabled interrupt or reset occurs.
	\item[\mnemonic{RSET}] Performs the reset sequence and jumps to the reset vector.
	\item[\mnemonic{ETFR}] Copies the saved return address and/or status register from a link register into \reg{a}
	      and/or \reg{r7}.
	\item[\mnemonic{ETTR}] Copies \reg{a} and/or \reg{r7} into a link register's saved return address and/or status
	      register.
\end{description}

\mnemonic{EXCP}, \mnemonic{RETE}, \mnemonic{WAIT}, \mnemonic{RSET}, \mnemonic{ETFR}, and \mnemonic{ETTR} are
meta-instructions that assemble to \mnemonic{COPI} (coprocessor call) forms addressed to the exception coprocessor
(coprocessor \texttt{0x1}); Section~\ref{sec:exception-registers-coprocessor-summary-forward} in this chapter lists
their opcodes, and Chapter~\ref{ch:coprocessor} gives the full syntax, operation, flags, and legal forms for each.

\section{Exception Types}

\subsection{Faults}

Faults are priority 7 exceptions raised by the CPU itself in response to error conditions or debug events. They use
vectors in the \texttt{0x01}--\texttt{0x0F} range. Unlike hardware and software exceptions, faults cannot be masked or
deferred: once a fault condition is raised, the fault is dispatched at the next instruction boundary regardless of the
current exception level. Two faults are conditional on whether they are raised at all: Segment Overflow requires
\texttt{SR.A}, and Instruction Trace requires \texttt{SR.T} or the \texttt{TRCE} input.

The T bit (TraceMode) in the status register is always cleared when entering any exception handler, including fault
handlers. This ensures that trace faults do not fire inside handlers. RETE restores the saved status register, so if T
was set before the exception, tracing resumes after RETE returns.

\subsubsection{Fault Categories}

Not all faults behave the same way with respect to the faulting instruction and the RETE return address. There are three
categories:

\begin{description}
	\item[Retryable faults] The operation that triggered the fault does not take effect. For most faults this means the
	      instruction is cancelled before it commits; Alignment Fault is the exception, since the fetch itself is what fails,
	      so there is no decoded instruction to cancel in the first place -- but the effect is the same either way. The link
	      register stores the address of the faulting instruction. RETE returns to that instruction so it can be retried, typically after
	      the fault handler has corrected the condition that caused the fault: remapping or paging in a device for a bus fault,
	      adjusting a permission for a bus protection fault, bringing an out-of-range address back into range for a segment
	      overflow, or emulating the missing operation for an invalid opcode fault. This is also why Privilege Violation is
	      retryable: the CPU guarantees that the instruction that triggered it never took effect, so a supervisor-mode handler
	      can safely inspect it, emulate the privileged effect if appropriate, and either retry it or advance past it under its
	      own control. Every fault in this category follows the same pattern: a handler is always free to advance past the
	      faulting instruction with \mnemonic{ETFR}/\mnemonic{ETTR} instead of retrying it. For Invalid Opcode Fault this is
	      usually the right choice -- simply retrying an unemulated invalid opcode will raise the same fault again.

	      \textit{Alignment Fault; Bus Fault; Bus Protection Fault; Invalid Opcode Fault; Privilege Violation Fault; Segment
		      Overflow Fault.}

	\item[Post-instruction faults] The instruction completes normally and its effects are committed before the fault fires. The
	      link register stores the address of the \textit{next} instruction. RETE advances to that instruction; the traced
	      instruction is not re-executed.

	      \textit{Instruction Trace Fault.}

	\item[Exception dispatch faults] Raised during the exception handling machinery itself, not during normal instruction
	      execution. The link register stores the program counter at the point the fault was detected within the exception
	      handler. Recovery may require careful state management; these faults often indicate a hardware misconfiguration or a
	      programming error in an exception handler.

	      \textit{Level Five Hardware Exception Conflict, Double Fault.}
\end{description}

\subsubsection{Fault Descriptions}

\begin{itemize}
	\item \textbf{Bus Fault} \textit{(Retryable)} — Raised when an external device asserts the bus error CPU pin.

	      This could happen, for example, if an unmapped address is presented by the CPU and another chip detects this and raises
	      an error. The instruction is cancelled before it takes effect, and the link register stores the address of the
	      \textit{faulting} instruction, so a handler implementing virtual memory can map in the missing page and return with
	      \mnemonic{RETE} to retry the same instruction.

	      Because ALU results and status-register updates are only committed during the write-back phase
	      (Chapter~\ref{ch:cpu-architecture}), which always runs after the memory-access phase where a bus fault is
	      detected, a bus fault can never leave a partial status-register or register update behind.

	\item \textbf{Bus Protection Fault} \textit{(Retryable)} — Raised when an external device asserts the bus
	      protection error CPU pin.

	      This indicates that the requested address and device are valid, but the operation is disallowed. Examples include
	      writing to ROM or accessing a protected memory region while the CPU is in protected mode. The instruction is cancelled
	      before it takes effect, and the link register stores the address of the faulting instruction, so a handler can adjust
	      the permission (if that is the correct response) and retry with \mnemonic{RETE}.

	\item \textbf{Alignment Fault} \textit{(Retryable)} — Raised when fetching instructions if the program counter
	      contains an odd word address.

	      This is to simplify fetching: the second word of an instruction is always at \texttt{pl | 0x1}, which ensures that
	      instructions can never overflow a segment boundary. The saved return address is itself odd; a fault handler must
	      correct \reg{p} (for example, by clearing bit 0) before returning with \mnemonic{RETE}, or the CPU raises the same
	      fault again immediately.

	\item \textbf{Segment Overflow Fault} \textit{(Retryable)} — Raised when a computed address would fall outside the
	      current segment.

	      Only raised if \texttt{SR.A} (Trap on Address Overflow) is set, since some programs intentionally rely on address
	      wrap-around. Also raised if the program counter overflows; the captured bus address and bus access type in the fault
	      metadata register help identify where the overflow occurred. The link register stores the address of the faulting
	      instruction, whether the overflow occurred in the program counter or in any other address computation, so a handler
	      can bring the address back into range and retry with \mnemonic{RETE}.

	\item \textbf{Invalid Opcode Fault} \textit{(Retryable)} — Raised when a coprocessor call targets a
	      non-existent coprocessor or requests an operation that the selected coprocessor does not implement.

	      Can be used for forward compatibility: programs written for a future coprocessor (for example, a floating-point unit)
	      will trap on older hardware, and the handler can emulate the operation in software. Because this fault is retryable,
	      the link register stores the address of the faulting instruction; a handler that has emulated the operation must use
	      \mnemonic{ETFR}/\mnemonic{ETTR} to redirect the return address to the instruction following the trapped call before
	      returning with \mnemonic{RETE} -- retrying the instruction unmodified will simply raise the same fault again.

	      Note that this fault is specific to coprocessor instructions. Other instruction encodings outside the documented
	      instruction set -- i.e. normal processing-unit encodings -- are reserved or architecturally undefined; they are not
	      required to raise an invalid-opcode fault.

	\item \textbf{Privilege Violation Fault} \textit{(Retryable)} — Raised when a privileged operation is attempted
	      while not in supervisor mode:
	      \begin{enumerate}
		      \item Directly writing to the high word of any address register
		      \item Directly writing to the status register
		      \item Address-register-pair write-back that would change a high address-register word
		      \item Executing a supervisor-only coprocessor operation
		      \item Triggering a software exception with a vector below \texttt{0x60}
	      \end{enumerate}

	      The illegal instruction never takes effect: it is cancelled before it writes anything, and the link register stores
	      its own address so \mnemonic{RETE} returns to it directly. This matters because a supervisor-mode handler needs to be
	      able to inspect the exact instruction that faulted, decide how to respond (deny it, emulate its effect, or promote the
	      caller), and then either retry it or use \mnemonic{ETFR}/\mnemonic{ETTR} to redirect the return address elsewhere.
	      None of that is possible if the illegal instruction had already been allowed to run.

	\item \textbf{Instruction Trace Fault} \textit{(Post-instruction)} — Raised after every instruction completes
	      when the \texttt{TraceMode} SR bit is set, or when the \texttt{TRCE} input is asserted at the instruction boundary
	      (see Appendix~\ref{appendix:external-interface}, Force Trace Mode Input).

	      Used for single-step debugging. Because the fault fires after the instruction commits, RETE returns to the instruction
	      \textit{following} the traced one. The T bit is cleared on entry to the fault handler (as with all exceptions), so the
	      handler itself is not traced. RETE restores the saved SR, which re-enables tracing unless the handler explicitly clears
	      T from the saved status register via \texttt{ETTR} before returning. \texttt{SR.T} is sampled once at the start of each
	      instruction, before that instruction can modify it, so a write to \texttt{SR.T} takes effect starting with the next
	      instruction boundary; exception-unit dispatches are never traced.

	\item \textbf{Level Five Hardware Exception Conflict} \textit{(Exception dispatch)} — Raised when a level five
	      hardware exception arrives while one is already being serviced.

	      The exception unit has no spare link register above level 6, so a second L5 cannot be handled normally. Rather than
	      silently dropping it (which could mask a hardware misconfiguration), the CPU raises this fault. Level 5 interrupts
	      should be treated as NMIs.

	\item \textbf{Double Fault} \textit{(Exception dispatch)} — Raised when a fault occurs while another fault is
	      already being handled.

	      When the CPU is already at fault level (priority 7) and another fault occurs, there is no additional banked link
	      register available. The fault link register (register 6) is overwritten with the new return address, losing the
	      original fault's return address. The fault metadata register (register 7) is also overwritten; it encodes a
	      double-fault flag along with both the current and original fault types.

	      The CPU jumps to vector \texttt{0x08}. This is typically a last-chance handler that dumps registers and resets. If
	      another fault occurs while the CPU is already at fault level, the CPU attempts to dispatch vector \texttt{0x08} again.
	      There is no limit to the number of times vector \texttt{0x08} is dispatched. A double fault handler that causes a
	      fault may cause an infinite loop without special handling from the programmer. If the double-fault vector fetch
	      itself fails, the CPU will reset. If two fault conditions are detected within the same instruction before either has
	      been dispatched, the resulting behavior is architecturally undefined.

\end{itemize}

\subsection{Hardware Exceptions (Interrupts)}

A hardware exception occurs when interrupt pins are signalled by external hardware. These are also referred to as
"interrupts". When a hardware exception occurs, the current instruction is completed, and control is transferred to
different vectors based on the interrupt "priority". The vectors for hardware exceptions are in the
\texttt{0x10}--\texttt{0x50} range.

The return address stored in the exception link register is the address of the instruction that would have executed
next, not the interrupted one -- so once the handler returns via \mnemonic{RETE}, the interrupted program resumes
normally as if nothing happened.

\subsection{User Exceptions (Traps)}

A user exception occurs when the program executes an \mnemonic{EXCP} instruction that signals the exception coprocessor
that it should trigger an exception. They are usually used by programs running in protected mode to switch to
supervisor mode in a controlled way. They are also referred to as "traps". The vectors for user exceptions are in the
\texttt{0x60}--\texttt{0xFF} range.

As with hardware exceptions, the return address stored in the link register is the instruction following the
\mnemonic{EXCP} call, so \mnemonic{RETE} resumes the calling program immediately after the trap.

Software exceptions are only accepted from normal execution. If an \mnemonic{EXCP} instruction is executed while any
exception or fault handler is already active, the software exception is ignored, is not queued, and does not fetch its
vector.

\section{Exception Registers}

The SIRC-1 exception coprocessor maintains several special-purpose registers for exception handling.

\subsection{Link Registers}

The exception unit uses banked link registers to store the processor state when an exception occurs. There are eight
link registers in total, indexed 0--7:

\begin{itemize}
	\item Register 0: Software exceptions (traps)
	\item Registers 1--5: Hardware exceptions (interrupts), one for each priority level
	\item Register 6: Faults
	\item Register 7: Fault metadata (stores bus address and bus access type information, see section
	      \ref{sec:fault-metadata-register-format} for the format)
\end{itemize}

Each link register stores three pieces of information:

\begin{itemize}
	\item \textbf{Return Address:} The 24-bit return address, held in a 32-bit field whose top 8 bits are always zero.
	      This is enforced by the same address masking applied whenever a full address is derived from an address register
	      pair (Figure~\ref{fig:addr-pair}), so even if the source address register's high word carries data in its own top
	      8 bits, that data is discarded before the address ever reaches the link register.
	\item \textbf{Return Status Register:} The 16-bit status register value at the time of the exception
	\item \textbf{Saved Exception Level:} The exception level that was active when the exception was taken;
	      \mnemonic{RETE} restores it into \texttt{current\_exception\_level}
\end{itemize}

\begin{figure}[H]
	\centering
	\begin{bytefield}[bitwidth=1.35em]{32}
		\bitheader{0,15,16,31} \\
		\bitbox{32}{Return Address} \\
		\bitbox{16}{Return Status Register} & \bitbox{8}{Saved Level} & \bitbox{8}{Reserved}
	\end{bytefield}
	\caption{Exception Link Register Saved-State Layout}
	\label{fig:exception-link-register-layout}
\end{figure}

The Saved Level and Reserved fields are shown for illustration only: \mnemonic{ETFR} and \mnemonic{ETTR} transfer only
the Return Address and the Return Status Register between a link register and \reg{a}/\reg{r7}. The saved exception
level is not accessible to software through these instructions.

When an exception occurs, the appropriate link register (determined by the exception priority level) is populated with
the current program counter and status register values. For faults in the retryable category (Alignment Fault, Bus
Fault, Bus Protection Fault, Invalid Opcode Fault, Privilege Violation Fault, and Segment Overflow Fault), the return
address is the address of the faulting instruction, so it can be retried. For regular exceptions (hardware and software
exceptions) and Instruction Trace Fault, the return address is the address after the current instruction; see the
"Fault Categories" section above for the exception-dispatch faults (Level Five Hardware Exception Conflict, Double
Fault), which instead store the program counter at the point the fault was detected.

\begin{table}[H]
	\centering
	\begin{tabular}{lll}
		\toprule
		\textbf{Register} & \textbf{Exception Source} & \textbf{Priority Level} \\
		\midrule
		0                 & Software exceptions       & 1                       \\
		1                 & Level 1 hardware          & 2                       \\
		2                 & Level 2 hardware          & 3                       \\
		3                 & Level 3 hardware          & 4                       \\
		4                 & Level 4 hardware          & 5                       \\
		5                 & Level 5 hardware          & 6                       \\
		6                 & Faults                    & 7                       \\
		7                 & Fault metadata            & n/a                     \\
		\bottomrule
	\end{tabular}
	\caption{Exception Link Register Assignment}
	\label{tab:exception-link-register-assignment}
\end{table}

\subsubsection{Link Register Preservation}

\textbf{Important:} Link registers can be overwritten when a fault occurs during the handling of another fault (a double fault).
In this case, the fault link register (register 6) and the fault metadata register (register 7) are overwritten.
This means the original return address is lost.

There is no way to fully recover from a double fault: the valid response is to declare the CPU state invalid, perform
some last-chance error handling (for example, dumping the registers), and reset the CPU.

A fault handler that wants a better chance of capturing useful diagnostics before that reset should still save the
link register contents early -- before doing anything else that might fault again -- using the \mnemonic{ETFR}
instruction, so the original fault's return address and status register can be included in whatever gets dumped.

\mnemonic{ETFR} only transfers the return address and status register of link registers 6 and 7; the saved exception
level is not accessible to software through it, nor can \mnemonic{ETTR} write it back. This is not just a missing
diagnostic field: \mnemonic{RETE} always restores \texttt{current\_exception\_level} directly from whatever value is
currently latched in the link register, and after a double fault that value is the fault level (7), not the level that
was active before the original fault. Even a handler that perfectly restores the return address and status register
and attempts \mnemonic{RETE} anyway would leave the CPU stuck at exception level 7 -- unable to service any hardware
interrupt or software exception, including Level 5, since nothing outranks the highest priority level. New faults
still dispatch normally regardless of the stuck level (faults bypass the priority check), so the CPU would keep
running rather than visibly failing; it would simply be deaf to every interrupt and trap for the rest of the program's
life. This is the underlying reason reset, not \mnemonic{RETE}, is the only valid response to a double fault.

The values can be stored on the stack but since the stack pointer may be corrupted, a more robust
implementation stores and loads from hardcoded memory addresses that are known to always be valid, instead of relying
on the stack pointer.

\subsubsection{Fault Metadata Register Format}
\label{sec:fault-metadata-register-format}

The fault metadata register (register 7) stores additional information about faults beyond just the return address. The
return\_status\_register field is used to encode fault-specific metadata in the following format:

\begin{figure}[H]
	\centering
	\begin{bytefield}[bitwidth=1.5em]{16}
		\bitheader{0,2,3,4,7,8,11,12,15} \\
		\bitbox{4}{Reserved} & \bitbox{4}{Original Fault} & \bitbox{4}{Current Fault} & \bitbox{1}{D} &
		\bitbox{3}{BAT}
	\end{bytefield}
	\caption{Fault Metadata Register Status Field Layout}
	\label{fig:fault-metadata-register-layout}
\end{figure}

\begin{itemize}
	\item \textbf{Bits 0--2:} Bus access type (\texttt{BAT0}--\texttt{BAT2}) captured when the fault was raised (3 bits)
	      \begin{itemize}
		      \item \texttt{0x0}: None
		      \item \texttt{0x1}: Instruction fetch
		      \item \texttt{0x2}: Data read
		      \item \texttt{0x3}: Data write
		      \item \texttt{0x4}: Exception vector fetch
		      \item \texttt{0x5}: DMA read burst
		      \item \texttt{0x6}: DMA write burst
		      \item \texttt{0x7}: Reserved
	      \end{itemize}
	\item \textbf{Bit 3:} Double fault flag (1 bit) — set to 1 if this is a double fault
	\item \textbf{Bits 4--7:} Current fault type (4 bits) — the type of fault that occurred
	\item \textbf{Bits 8--11:} Original fault type (4 bits) — only meaningful for double faults, indicates the fault type
	      that was being handled when the double fault occurred
	\item \textbf{Bits 12--15:} Unused (reserved for future use)
\end{itemize}

Fault types use the following values:

\begin{itemize}
	\item \textbf{\texttt{0x01}:} Bus fault
	\item \textbf{\texttt{0x02}:} Alignment fault
	\item \textbf{\texttt{0x03}:} Segment overflow fault
	\item \textbf{\texttt{0x04}:} Invalid opcode fault
	\item \textbf{\texttt{0x05}:} Privilege violation fault
	\item \textbf{\texttt{0x06}:} Instruction trace fault
	\item \textbf{\texttt{0x07}:} Level five hardware exception conflict
	\item \textbf{\texttt{0x08}:} Double fault
	\item \textbf{\texttt{0x09}:} Bus protection fault
\end{itemize}

The return\_address field of register 7 stores the bus address that was being accessed when the fault occurred (for
faults that involve memory access). For coprocessor invalid-opcode faults that occur during coprocessor dispatch before
bus access, the low word of this field stores the 16-bit coprocessor command that could not be dispatched, and the high
word is zero. This lets an invalid-opcode handler emulate missing optional coprocessor operations.

Exception handlers can use the \mnemonic{ETFR} instruction to read the fault metadata register and decode this
information to determine how to handle the fault appropriately.

\subsection{Cause Register}

The cause register is a 16-bit internal register used to communicate between different CPU units and coordinate
exception handling. It has the following structure:

\begin{itemize}
	\item \textbf{Bits 15--12 (First Nibble):} Coprocessor ID (\texttt{0x1} for exception unit, \texttt{0x0} for processing
	      unit)
	\item \textbf{Bits 11--8 (Second Nibble):} Coprocessor opcode (the operation to execute)
	\item \textbf{Bits 7--0 (Third and Fourth Nibbles):} Vector or parameter value
\end{itemize}

The cause register is automatically populated by the exception unit when a pending exception needs to be serviced. The
priority of exceptions is determined by examining the first nibble of the vector ID, calculated as 7 minus the value of
the first nibble (for example, \texttt{0x00} is priority 7, \texttt{0x40} is priority 3, and \texttt{0x60} and above
are all priority 1).

A cause value of \texttt{0x0000} -- the value left in place at reset and whenever no exception is pending -- selects
coprocessor ID \texttt{0x0}, the processing unit, so the CPU simply continues executing the program normally in the
absence of any exception activity.

\subsection{Hardware Interrupt Enable Bits}
\label{sec:hardware-interrupt-enable-bits}

The status register contains 4 individual enable bits (bits 9--12 in the privileged byte) that control whether each
maskable hardware interrupt line can trigger an exception:

\begin{itemize}
	\item Bit 9 (E1): Hardware Interrupt Line 1 (Level 1, lowest maskable priority)
	\item Bit 10 (E2): Hardware Interrupt Line 2 (Level 2)
	\item Bit 11 (E3): Hardware Interrupt Line 3 (Level 3)
	\item Bit 12 (E4): Hardware Interrupt Line 4 (Level 4, highest maskable priority)
\end{itemize}

\textbf{Note:} Hardware Interrupt Line 5 (Level 5) is a Non-Maskable Interrupt (NMI) and cannot be disabled.
It will always trigger an exception when asserted, regardless of the enable bit settings. Bit 13 is used for the
Exception Active flag instead.

When an interrupt line is disabled (bit = 0), interrupts on that line are completely ignored and not queued. When
enabled (bit = 1), the interrupt can trigger an exception if its priority is higher than the currently executing
exception level.

\subsection{Pending Hardware Interrupts}

Hardware interrupt lines 1--4 are level-sensitive inputs; Level 5 is edge-triggered and is described separately under
"Level 5 NMI-like Behavior" below. When an enabled level-sensitive interrupt line is observed asserted, the
corresponding bit is set in the pending hardware-exceptions register. Disabled maskable interrupt lines are ignored at
this sampling point and are not latched for later service.

For the level-sensitive lines, the pending register stores one bit per hardware interrupt line, so repeated assertions
of the same line coalesce into a single pending interrupt. If the external pin remains asserted after the interrupt is
serviced, it may be sampled again and become pending again.

At an instruction boundary, the exception unit selects the highest-priority pending hardware interrupt whose priority
is higher than the current exception level. The selected bit is cleared as the handler is dispatched. Pending
lower-priority interrupts remain latched and are considered again after the current handler returns with
\mnemonic{RETE}.

Faults have priority over all hardware interrupts. If a fault is pending at the same instruction boundary as one or
more hardware interrupts, the fault is dispatched first and the hardware interrupts remain pending. This includes
instruction trace faults: when trace mode raises a post-instruction fault, that trace fault is serviced before pending
hardware interrupts at the next boundary.

\textbf{Hierarchical Exception Handling:}

Exceptions in SIRC-1 are strictly hierarchical. A hardware interrupt can only interrupt the current execution if its
priority level (exception level) is \textit{higher} than the currently active exception level. This prevents:

\begin{itemize}
	\item An interrupt handler from being interrupted by its own interrupt (same level)
	\item Lower priority interrupts from interrupting higher priority handlers
	\item Software exceptions from running during any hardware interrupt (would corrupt link register 0)
\end{itemize}

The exception unit maintains a \texttt{current\_exception\_level} field that tracks the active exception level (0--7).
When not handling any exception, this is 0. When handling an exception at level N, only exceptions at level > N can
interrupt.

\textbf{Level 5 NMI-like Behavior:}

Level 5 hardware exceptions behave similarly to a Non-Maskable Interrupt (NMI). Unlike interrupt levels 1--4, which are
level-sensitive (sampled for as long as the pin is asserted), Level 5 is edge-triggered: it fires once on assertion
regardless of how long the pin is held, and must be released and reasserted to fire again. This avoids spurious Level
Five Hardware Exception Conflict faults from a line that simply stays asserted. Level 5 cannot be disabled via the
interrupt enable bits in the status register, and will preempt any lower-priority exception handler. However, Level 5
is \textit{not} fully non-maskable. Because faults carry priority 7 (higher than Level 5's priority 6), a L5 interrupt
that arrives while a fault is being handled is \textit{queued} in the pending hardware exceptions register rather than
serviced immediately.

Once the fault handler returns via \texttt{RETE} and the exception level drops back below 6, the pending L5 interrupt
fires normally.

This means Level 5 response latency is not deterministic when a fault is in progress. For timing-critical use of the
NMI line, fault handlers should be kept as short as possible.

If a second L5 interrupt arrives \textit{while the CPU is already executing an L5 handler} (exception level 6), it is
not queued. It will raise a Level Five Hardware Exception Conflict fault instead. See the fault descriptions above.

These two rules compose in a non-obvious way when a fault is raised \textit{inside} an L5 handler while a second L5
interrupt is pending. Table~\ref{tab:l5-fault-timeline} walks through the sequence step by step.

\begin{table}[H]
	\centering
	\begin{tabularx}{\textwidth}{c>{\raggedright\arraybackslash}p{4cm}>{\raggedright\arraybackslash}p{2.3cm}>{\raggedright\arraybackslash}X}
		\toprule
		\textbf{Step} & \textbf{Event}                                                                    & \textbf{Exception Level} & \textbf{Link Register Used}                                                        \\
		\midrule
		1             & CPU is executing inside the L5 handler.                                          & 6                        & Link register 5 (holds the L5 handler's return state)                             \\
		2             & A fault condition is raised inside the L5 handler. Priority 7 exceeds the current level (6), so the fault preempts it. & 6 $\rightarrow$ 7 & Link register 6 receives the L5 handler's return address, status register, and saved level (6) \\
		3             & A second L5 interrupt is asserted while the fault handler runs. Priority 6 does not exceed the current level (7), so it is not serviced. & 7 & Not dispatched; latched as pending in the pending hardware-exceptions register \\
		4             & The fault handler executes \mnemonic{RETE}, restoring the return address, status register, and saved level (6) from link register 6. & 7 $\rightarrow$ 6 & Link register 6 (restored, not overwritten) \\
		5             & Back at level 6, the exception unit reconsiders the pending L5 interrupt. Its priority (6) is not \textit{higher} than the current level (6), so it is not dispatched normally. & 6 & --- \\
		6             & The CPU raises a Level Five Hardware Exception Conflict fault instead, because the pending L5 collides with the still-active L5 handler. & 6 $\rightarrow$ 7 & Link register 6 receives the dispatch point as the return address \\
		\bottomrule
	\end{tabularx}
	\caption{Fault-Inside-L5-Handler Timeline}
	\label{tab:l5-fault-timeline}
\end{table}

In short: the fault (priority 7) blocks the second L5 (priority 6) from firing while the fault handler runs. Once
\mnemonic{RETE} returns to level 6, the queued L5 collides with the still-active outer L5 handler, and a Level Five
Hardware Exception Conflict fault is raised then, not during the inner fault handler.

\section{Exception Processing}

\subsection{Exception Priority}

Exceptions are prioritized as follows (highest to lowest):

\begin{enumerate}
	\item Priority 7: Faults
	\item Priority 6: Level 5 hardware exceptions (NMI-like)
	\item Priority 5: Level 4 hardware exceptions
	\item Priority 4: Level 3 hardware exceptions
	\item Priority 3: Level 2 hardware exceptions
	\item Priority 2: Level 1 hardware exceptions
	\item Priority 1: Software exceptions (traps)
\end{enumerate}

\subsection{Exception Handling Flow}

When an exception occurs, the following steps are executed:

\begin{enumerate}
	\item For hardware exceptions, the exception unit checks if the interrupt line is enabled in the status register. If
	      disabled, the interrupt is ignored and not queued.

	\item The exception unit checks if the exception priority level is higher than the current active exception level. If not,
	      the exception is deferred until the current handler completes.

	\item For level 5 hardware exceptions being raised when one is already being handled (both have exception level 6), a Level
	      Five Hardware Exception Conflict fault is raised instead.

	\item The vector-table address is computed by multiplying the vector ID by 2, since each vector is a 24-bit address stored as
	      two 16-bit words.

	\item The exception unit reads the high word and then the low word of the vector target address from the vector table.

	\item The current program counter, status register, and current exception level are stored in the appropriate link register
	      based on the exception priority level.

	\item The Protected Mode bit in the status register is cleared, switching the CPU to supervisor mode.

	\item The current exception level is updated to the level of the exception being serviced.

	\item The program counter is set to the fetched vector target address, transferring control to the exception handler.
\end{enumerate}

\subsection{Exception Entry Side Effects}
\label{sec:exception-entry-side-effects}

For pending hardware interrupts, priority and mask checks occur before the exception-vector fetch. A hardware interrupt
that is disabled or not high enough priority to preempt the current handler does not fetch its vector; if it is already
latched as pending, it remains pending for later service according to the pending-interrupt rules above.

Software exceptions are also screened before the exception-vector fetch. A software exception is accepted only when no
exception handler is active.

After an exception source is selected for service, the exception unit fetches the high word and low word of the handler
target address from the vector table. Architectural entry side effects commit only after the vector target has been
fetched and the exception is accepted for service.

When entry commits, the exception unit performs the following state updates as one architectural action:

\begin{itemize}
	\item Writes the selected link register with the current program counter, current status register, and current exception
	      level.
	\item Clears \texttt{SR.P}, entering supervisor mode.
	\item Clears \texttt{SR.T}, preventing trace faults from firing inside the handler.
	\item Sets \texttt{SR.EA}, indicating that an exception handler is active.
	\item Preserves the condition flags, reserved lower status bits, interrupt-enable bits, and \texttt{SR.A} in the live status
	      register.
	\item Updates the current exception level to the level of the exception being serviced.
	\item Loads the program counter from the fetched vector target address.
\end{itemize}

The saved status register in the link register is the pre-entry value, before \texttt{SR.P}, \texttt{SR.T}, or
\texttt{SR.EA} are modified. \mnemonic{RETE} restores that saved status register, so a handler can inspect or edit the
saved value with \mnemonic{ETFR} and \mnemonic{ETTR} before returning.

If an exception is not accepted because its priority is not higher than the current exception level, none of these
entry side effects occur. Pending hardware interrupts remain subject to the priority and pending-interrupt rules
described above.

If a bus or protection fault occurs during an exception-vector fetch before entry commits, the original exception does
not enter. The CPU raises the corresponding fault with \texttt{BAT0}--\texttt{BAT2} set to
\texttt{ExceptionVectorFetch}. If a fault-vector fetch itself fails, the CPU escalates to double fault. If the
double-fault vector fetch fails, the CPU will reset.

\textbf{Returning from Exceptions:} see Section~\ref{sec:return-from-exception} for how \mnemonic{RETE} restores this
state.

\subsection{Vector Table}

The exception vectors are stored starting at address \addr{0x0}. Each vector is a 24-bit target address stored as two
consecutive 16-bit words in the address format described in Section~\ref{sec:exception-vector-storage}. The high word
is stored first at \texttt{vector * 2}, followed by the low word at \texttt{vector * 2 + 1}. The vector table has 256
entries. Vectors \texttt{0x00}--\texttt{0x5F} (96 entries) are reserved for the CPU; vectors
\texttt{0x60}--\texttt{0xFF} (160 entries) are user trap vectors. The vector table layout is:

\begin{itemize}
	\item \textbf{\texttt{0x00}:} Reset vector
	\item \textbf{\texttt{0x01}:} Bus fault
	\item \textbf{\texttt{0x02}:} Alignment fault
	\item \textbf{\texttt{0x03}:} Segment overflow fault
	\item \textbf{\texttt{0x04}:} Invalid opcode fault
	\item \textbf{\texttt{0x05}:} Privilege violation fault
	\item \textbf{\texttt{0x06}:} Instruction trace fault
	\item \textbf{\texttt{0x07}:} Level five hardware exception conflict
	\item \textbf{\texttt{0x08}:} Double fault
	\item \textbf{\texttt{0x09}:} Bus protection fault
	\item \textbf{\texttt{0x0A}--\texttt{0x0F}:} Reserved for future privileged exceptions
	\item \textbf{\texttt{0x10}:} Level 5 hardware exception vector
	\item \textbf{\texttt{0x20}:} Level 4 hardware exception vector
	\item \textbf{\texttt{0x30}:} Level 3 hardware exception vector
	\item \textbf{\texttt{0x40}:} Level 2 hardware exception vector
	\item \textbf{\texttt{0x50}:} Level 1 hardware exception vector
	\item \textbf{\texttt{0x11}--\texttt{0x5F} (excluding \texttt{0x20}, \texttt{0x30}, \texttt{0x40}, \texttt{0x50}):} Reserved
	\item \textbf{\texttt{0x60}--\texttt{0xFF}:} User exception vectors (160 user-accessible trap vectors)
\end{itemize}

\subsection{Exception Quick Reference}

\begin{table}[H]
	\centering
	\scriptsize
	\begin{tabularx}{\textwidth}{>{\raggedright\arraybackslash}p{1.2cm}>{\raggedright\arraybackslash}p{1.5cm}>{\raggedright\arraybackslash}p{1.1cm}>{\raggedright\arraybackslash}p{1.9cm}>{\raggedright\arraybackslash}p{1.7cm}>{\raggedright\arraybackslash}X>{\raggedright\arraybackslash}p{1.5cm}>{\raggedright\arraybackslash}p{1.2cm}}
		\toprule
		\textbf{Vector}              & \textbf{Address}             & \textbf{Priority} & \textbf{Source}  & \textbf{Maskability} & \textbf{Return Address}                   & \textbf{Metadata} & \textbf{P Mode} \\
		\midrule
		\texttt{0x00}                & \addr{0x0000}--\addr{0x0001} & n/a               & Reset            & Not maskable         & n/a                                       & None              & Yes             \\
		\texttt{0x01}                & \addr{0x0002}--\addr{0x0003} & 7                 & Bus fault        & Not maskable         & Next instruction                          & Fault metadata    & Yes             \\
		\texttt{0x02}                & \addr{0x0004}--\addr{0x0005} & 7                 & Alignment fault  & Not maskable         & Faulting instruction                      & Fault metadata    & Yes             \\
		\texttt{0x03}                & \addr{0x0006}--\addr{0x0007} & 7                 & Segment overflow & \mbox{SR.A gated}    & Faulting instruction                      & Fault metadata    & Yes             \\
		\texttt{0x04}                & \addr{0x0008}--\addr{0x0009} & 7                 & Invalid opcode   & Not maskable         & Faulting instruction                      & Fault metadata    & Yes             \\
		\texttt{0x05}                & \addr{0x000A}--\addr{0x000B} & 7                 & Privilege fault  & Not maskable         & Next instruction                          & Fault metadata    & Yes             \\
		\texttt{0x06}                & \addr{0x000C}--\addr{0x000D} & 7                 & Trace fault      & SR.T or TRCE gated   & Next instruction                          & Fault metadata    & Yes             \\
		\texttt{0x07}                & \addr{0x000E}--\addr{0x000F} & 7                 & L5 conflict      & Not maskable         & Dispatch point                            & Fault metadata    & Yes             \\
		\texttt{0x08}                & \addr{0x0010}--\addr{0x0011} & 7                 & Double fault     & Not maskable         & Dispatch point                            & Fault metadata    & Yes             \\
		\texttt{0x09}                & \addr{0x0012}--\addr{0x0013} & 7                 & Bus protection   & Not maskable         & Next instruction                          & Fault metadata    & Yes             \\
		\texttt{0x0A}--\texttt{0x0F} & \addr{0x0014}--\addr{0x001F} & 7                 & Reserved faults  & Reserved             & Reserved                                  & Reserved          & Reserved        \\
		\texttt{0x10}                & \addr{0x0020}--\addr{0x0021} & 6                 & Level 5 hardware & Not maskable         & Next instruction                          & Link register 5   & Yes             \\
		\texttt{0x20}                & \addr{0x0040}--\addr{0x0041} & 5                 & Level 4 hardware & \mbox{SR.E4}         & Next instruction                          & Link register 4   & Yes             \\
		\texttt{0x30}                & \addr{0x0060}--\addr{0x0061} & 4                 & Level 3 hardware & \mbox{SR.E3}         & Next instruction                          & Link register 3   & Yes             \\
		\texttt{0x40}                & \addr{0x0080}--\addr{0x0081} & 3                 & Level 2 hardware & \mbox{SR.E2}         & Next instruction                          & Link register 2   & Yes             \\
		\texttt{0x50}                & \addr{0x00A0}--\addr{0x00A1} & 2                 & Level 1 hardware & \mbox{SR.E1}         & Next instruction                          & Link register 1   & Yes             \\
		\texttt{0x11}--\texttt{0x5F} & \addr{0x0022}--\addr{0x00BF} & n/a               & Reserved         & Reserved             & Reserved                                  & Reserved          & Reserved        \\
		\texttt{0x60}--\texttt{0xFF} & \addr{0x00C0}--\addr{0x01FF} & 1                 & Software trap    & Priority gated       & Next instruction                          & Link register 0   & Yes             \\
		\bottomrule
	\end{tabularx}
	\caption{Exception Quick Reference}
	\label{tab:exception-quick-reference}
\end{table}

\section{Reset and Startup}

Reset can be initiated externally by asserting the \texttt{RSTI} pin, or in software by executing the privileged
\mnemonic{RSET} meta-instruction. Both paths converge on the same reset-vector sequence: the CPU holds reset output
active, then the exception unit fetches vector \texttt{0x00} and loads the fetched target address into the program
counter.

\subsection{Reset State}

\begin{table}[H]
	\centering
	\begin{tabularx}{\textwidth}{l>{\raggedright\arraybackslash}X}
		\toprule
		\textbf{State}                         & \textbf{Reset Behavior}                                                                              \\
		\midrule
		Instruction phase                      & Set to phase 0.                                                                                      \\
		Pending bus request                    & Cleared; an in-progress stalled bus request is abandoned.                                            \\
		Halt state                             & Cleared.                                                                                             \\
		Wait-for-exception state               & Cleared.                                                                                             \\
		Current exception level                & Cleared to 0. The CPU leaves reset outside any exception handler.                                    \\
		Status register                        & Cleared to \texttt{0x0000} when the reset operation is executed by the exception unit.               \\
		Program counter                        & Loaded from reset vector \texttt{0x00}.                                                              \\
		General registers                      & Not architecturally cleared by reset. Software should treat their contents as undefined after reset. \\
		Address register pairs                 & Not architecturally cleared by reset.                                                                \\
		Exception link registers               & Not architecturally cleared by reset.                                                                \\
		Pending hardware exceptions and faults & Cleared. External interrupt pins that remain asserted after reset may be
		sampled again according to the normal interrupt rules.                                                                                        \\
		Pending coprocessor command            & Set internally to the reset operation until the reset vector sequence completes.                     \\
		\bottomrule
	\end{tabularx}
	\caption{Reset State}
	\label{tab:reset-state}
\end{table}

Because the status register is cleared to \texttt{0x0000} as part of the reset sequence, hardware interrupt lines 1--4
are masked immediately after reset -- their enable bits, \texttt{SR.E1}--\texttt{SR.E4}
(Section~\ref{sec:hardware-interrupt-enable-bits}), are all cleared to 0. The Level 5 (NMI) line is architecturally
non-maskable and remains enabled regardless of \reg{sr}'s contents. Note also that the status register is not cleared
the instant reset is requested: it is cleared during the reset-vector fetch sequence (Section~\ref{sec:reset-vector-fetch}),
once that sequence reaches its execute phase.

\subsection{Reset Output Hold}
\label{sec:reset-output-hold}

While reset is active, the CPU asserts \texttt{RSTO} and performs no bus cycles. This gives external devices sharing
the reset line time to complete their own reset sequence before the CPU resumes.

For hardware reset, asserting \texttt{RSTI} immediately aborts any current instruction, exception handler, exception
dispatch, or pending bus wait, then asserts \texttt{RSTO}. If \texttt{RSTI} remains asserted, \texttt{RSTO} remains
asserted and the CPU does not advance. After \texttt{RSTI} is released, the reset hold completes before the CPU
resumes. The reset assertion cycle counts as the first reset-output cycle; after release, five additional countdown
cycles occur before the exception unit starts the reset-vector fetch.

For software reset, \mnemonic{RSET} completes as a normal coprocessor meta-instruction. When its write-back phase
commits the reset command, the CPU requests reset and \texttt{RSTO} is asserted for six cycles before the reset-vector
fetch begins.

\subsection{Reset Vector Fetch}
\label{sec:reset-vector-fetch}

The reset vector is vector \texttt{0x00} in the exception vector table. After the reset-output hold, the exception unit
performs an exception-vector fetch with bus access type \texttt{ExceptionVectorFetch}:

\begin{enumerate}
	\item Read the high word of the reset target address from \addr{0x0000}.
	\item Read the low word of the reset target address from \addr{0x0001}.
	\item Clear the status register to \texttt{0x0000}.
	\item Load the fetched target address into \reg{p}.
\end{enumerate}

These four steps are not one clock cycle each. Steps 1 and 2 each occupy their own fetch cycle, an instruction-decode
cycle follows between step 2 and step 3 (not itemized above), and steps 3 and 4 then happen together in a single
execute cycle. The reset-vector fetch therefore takes four clock cycles in total, separate from and following the
reset-output hold described in Section~\ref{sec:reset-output-hold}.

Step 3's status-register clear is listed here, rather than only in the Reset State table
(Table~\ref{tab:reset-state}), because it marks the specific point in the pipeline where the clear actually happens --
useful for anyone modeling cycle-accurate timing. Table~\ref{tab:reset-state} instead documents the eventual, settled
reset state; no other register is touched during the reset-vector fetch itself, so it is the only state change
covered both here and there.

If a bus or protection fault occurs during the reset-vector fetch, the reset target is not loaded. The CPU raises the
corresponding fault with \texttt{BAT0}--\texttt{BAT2} set to \texttt{ExceptionVectorFetch}. From that point the normal
exception-vector fetch fault rules apply: a fault-vector fetch failure escalates to double fault, and a double-fault
vector fetch failure will reset the CPU.

\section{Return from Exception}
\label{sec:return-from-exception}

To return from an exception handler, the \mnemonic{RETE} (Return from Exception) instruction restores the saved state
from the current link register:

\begin{enumerate}
	\item The return address is loaded from the link register into the program counter
	\item The saved status register value is restored (including interrupt enable bits)
	\item The saved exception level is restored to \texttt{current\_exception\_level}
	\item If the restored exception level is 0, \texttt{SR.EA} is cleared; otherwise the restored value of \texttt{SR.EA} stands
	\item Execution resumes at the return address
\end{enumerate}

After RETE, any pending interrupts that are now eligible to run (higher priority than the restored exception level and
enabled) will be serviced before the next instruction executes.

Executing \mnemonic{RETE} while the current exception level is 0 is architecturally undefined.

For faults in the retryable category (Alignment Fault, Bus Fault, Bus Protection Fault, Invalid Opcode Fault, Privilege
Violation Fault, and Segment Overflow Fault), \mnemonic{RETE} returns to the faulting instruction, allowing it to be
retried. For all other faults and for regular exceptions, \mnemonic{RETE} returns to the instruction following the one
that was executing when the exception occurred.

\subsection{Accessing Link Registers}

The exception coprocessor provides instructions to transfer data to and from the link registers:

\begin{itemize}
	\item \textbf{\mnemonic{ETFR} (Exception Transfer From Register):} Copies the return address and/or status register from a link register to the address register and/or \reg{r7}
	\item \textbf{\mnemonic{ETTR} (Exception Transfer To Register):} Copies the address register and/or \reg{r7} to a link register's return address and/or status register
\end{itemize}

These instructions allow exception handlers to examine or modify the saved processor state before returning, enabling
features like system calls, context switching, and debuggers.

\section{Exception Coprocessor Instructions}
\label{sec:exception-registers-coprocessor-summary-forward}

The exception coprocessor (coprocessor ID \texttt{0x1}) provides several specialized instructions for exception
handling and system control.

\begin{table}[H]
	\centering
	\begin{tabularx}{\textwidth}{lll>{\raggedright\arraybackslash}X}
		\toprule
		\textbf{Instruction} & \textbf{Opcode} & \textbf{Privilege} & \textbf{Purpose}                                            \\
		\midrule
		\mnemonic{EXCP}      & \texttt{0x1}    & User               & Trigger a software exception (trap)                         \\
		\mnemonic{WAIT}      & \texttt{0x9}    & Supervisor         & Enter wait state until an enabled interrupt or reset occurs \\
		\mnemonic{RETE}      & \texttt{0xA}    & Supervisor         & Return from an exception handler                            \\
		\mnemonic{RSET}      & \texttt{0xB}    & Supervisor         & Perform reset sequence and jump to reset vector             \\
		\mnemonic{ETFR}      & \texttt{0xC}    & Supervisor         & Copy saved exception state from a link register             \\
		\mnemonic{ETTR}      & \texttt{0xD}    & Supervisor         & Write CPU register state back to a link register            \\
		\bottomrule
	\end{tabularx}
	\caption{Exception Coprocessor Instruction Summary}
	\label{tab:exception-coprocessor-summary}
\end{table}

Exception coprocessor opcodes \texttt{0x2}--\texttt{0x8} are reserved; issuing one is architecturally undefined.

For full syntax, operation details, and worked examples, see Chapter~\ref{ch:coprocessor}.

\subsection{Internal Instructions}

The following opcodes are used internally by the exception unit and are not directly accessible via user instructions:

\begin{itemize}
	\item \textbf{Fault (\texttt{0xE}):} Automatically invoked when a fault condition is detected
	\item \textbf{HardwareException (\texttt{0xF}):} Automatically invoked when a hardware interrupt is signaled
	\item \textbf{None (\texttt{0x0}):} No operation, used when no exception is pending
\end{itemize}

These internal operations follow the standard exception handling flow described in the Exception Processing section.
Issuing opcode \texttt{0xE} or \texttt{0xF} directly, for example via a supervisor-mode \mnemonic{COPI} instruction, is
architecturally undefined.
